\documentclass[11pt]{article}
\usepackage[utf8]{inputenc}
\usepackage{ctex}
\usepackage{amsmath, amssymb, amsthm}
\usepackage{geometry}
\geometry{a4paper, margin=1in}

\input{../preamable/preamable.tex}

\declaretheorem[style=thmgreenbox, name=定义]{cndefinition}
\declaretheorem[style=thmredbox, name=定理]{cntheorem}
\declaretheorem[style=thmbluebox, name=性质]{cnproperty}
\declaretheorem[style=thmredbox, name=引理]{cnlemma}
\declaretheorem[style=thmredbox, numbered=no, name=推论]{cncorollary}
\declaretheorem[style=thmbluebox, numbered=no, name=例]{cnexample}
\title{近世代数笔记 (Lec 07)：环与域的基本概念}
\author{Raysance}
\date{}

\begin{document}

\maketitle

\section{环的定义与基本概念}

\begin{cndefinition}[环 (Ring)]
非空集合 $R$，连同两个二元运算（通常记为加法 $+$ 和乘法 $\cdot$，即 $R \times R \to R$），如果满足以下条件，则称 $(R, +, \cdot)$ 为一个\textbf{环}：
\begin{enumerate}
    \item $(R, +)$ 是阿贝尔群（Abelian Group）。这意味着加法满足交换律和结合律，存在加法单位元（称为零元，记为 $0$），且每个元素有加法逆元（记为 $-\alpha$）。
    \item $(R, \cdot)$ 是半群（Semigroup）。这意味着乘法满足结合律。
    \item 乘法对加法满足分配律。即对于任意 $\alpha, \beta, \gamma \in R$，有：
    \begin{itemize}
        \item 左分配律：$\alpha(\beta + \gamma) = \alpha\beta + \alpha\gamma$
        \item 右分配律：$(\beta + \gamma)\alpha = \beta\alpha + \gamma\alpha$
    \end{itemize}
\end{enumerate}
\end{cndefinition}

直观理解：环是整数集合的抽象。在环中，我们可以像操作整数一样自由地进行加法、减法和乘法（满足结合律和分配律），但除法（或乘法逆元）不一定存在，乘法也不一定满足交换律。

\begin{cndefinition}[含幺环与交换环]
\begin{enumerate}
    \item 如果环 $R$ 中存在一个元素 $e \in R$，使得对于所有的 $\alpha \in R$，都有 $e\alpha = \alpha e = \alpha$，那么 $e$ 称为乘法幺元（通常记为 $1$），此时称 $R$ 为\textbf{含幺环} (Ring with identity)。
    \item 如果环 $R$ 的乘法满足交换律（即对于任意 $\alpha, \beta \in R$，$\alpha\beta = \beta\alpha$），则称 $R$ 为\textbf{交换环} (Commutative ring)。
\end{enumerate}
\end{cndefinition}

\section{环的基本性质}

\begin{cnproperty}[环的基本运算性质]
在任意环 $R$ 中，对于任意 $\alpha, \beta \in R$，下列性质成立：
\begin{enumerate}
    \item $0\alpha = \alpha 0 = 0$
    \item $(-\alpha)\beta = \alpha(-\beta) = -(\alpha\beta)$
    \item $(\sum \alpha_i)\beta = \sum (\alpha_i\beta)$
    \item 对于整数 $n$，$(n\alpha)\beta = \alpha(n\beta) = n(\alpha\beta)$
\end{enumerate}
\end{cnproperty}
\begin{proof}
我们可以利用分配律从阿贝尔群的性质推导这些结论：
\begin{enumerate}
    \item 因为 $0+0=0$，所以 $0\alpha = (0+0)\alpha = 0\alpha + 0\alpha$。在加法群中消去 $0\alpha$，得到 $0\alpha = 0$。同理可证 $\alpha0=0$。
    \item $0 = 0\beta = (\alpha + (-\alpha))\beta = \alpha\beta + (-\alpha)\beta$。由加法逆元的唯一性，$(-\alpha)\beta$ 就是 $\alpha\beta$ 的加法逆元，即 $(-\alpha)\beta = -(\alpha\beta)$。同理可证 $\alpha(-\beta) = -(\alpha\beta)$。
    \item 与 4 同理，利用分配律的直接扩展。其中 $n\alpha$ 是指 $n$ 个 $\alpha$ 连加。
\end{enumerate}
\end{proof}

\section{零因子、可逆元与特殊环}

\begin{cndefinition}[零因子 (Zero Divisor)]
在环 $R$ 中，若对于非零元素 $\alpha \in R$，存在非零元素 $\beta \in R$ 使得 $\alpha\beta = 0$，则称 $\alpha$ 为\textbf{左零因子}，$\beta$ 称为\textbf{右零因子}。如果 $\alpha$ 同时是左零因子和右零因子，则称其为\textbf{零因子}。
\end{cndefinition}

直观理解：零因子是可能导致“非零元素相乘为零”的元素。整数中没有零因子，但在模合数 $n$ 的剩余类环中（如 $\mathbb{Z}_6$ 中的 $2 \times 3 = 0$），存在大量零因子。

\begin{cndefinition}[可逆元 \& 单位 (Unit)]
设 $R$ 为含幺环（具有幺元 $1$）。对于元素 $u \neq 0, u \in R$，若存在 $v \neq 0, v \in R$ 使得 $uv = 1$，则称 $u$ 为\textbf{右可逆}，$v$ 称为 $u$ 的\textbf{右逆}（相应的 $u$ 是 $v$ 的左逆）。
若 $u$ 同时左可逆且右可逆，则称 $u$ 为 $R$ 的\textbf{可逆元}（单位，Unit）。
\end{cndefinition}

\begin{cnproperty}[有关可逆元的评论与性质]
\begin{enumerate}
    \item 若某元素 $u$ 的左逆和右逆都存在，则它的左逆和右逆必定相等。
    \item 对于交换环，左逆天然等于右逆。
    \item 如果只有单侧逆存在，则逆元不一定唯一。
\end{enumerate}
\end{cnproperty}
\begin{proof}
对于第一点，设 $v$ 是 $u$ 的左逆（即 $vu=1$），$v'$ 是 $u$ 的右逆（即 $uv'=1$），则有：
$$v = v \cdot 1 = v(uv') = (vu)v' = 1 \cdot v' = v'$$
这就证明了只要两侧逆存在必然相等，且此时可逆元 $u$ 的逆是唯一的（通常记为 $u^{-1}$）。
\end{proof}

\begin{cnproperty}
设 $R$ 为含幺环，$U(R)$ 为 $R$ 中所有可逆元（单位）构成的集合。则 $U(R)$ 关于 $R$ 的乘法构成一个群。特别地：两个可逆元的乘积仍是一个可逆元。
\end{cnproperty}
\begin{proof}
对任意 $u_1, u_2 \in U(R)$，有 $(u_1 u_2)(u_2^{-1} u_1^{-1}) = u_1 (u_2 u_2^{-1}) u_1^{-1} = u_1 \cdot 1 \cdot u_1^{-1} = 1$，同理 $(u_2^{-1} u_1^{-1})(u_1 u_2) = 1$。说明 $u_1 u_2 \in U(R)$且 $(u_1 u_2)^{-1} = u_2^{-1} u_1^{-1}$。
\end{proof}

\begin{cndefinition}[整环与域]
\begin{enumerate}
    \item \textbf{整环 (Integral Domain)}: 一个含幺交换环，且没有零因子。
    \item \textbf{零环 (Zero Ring)}: 仅含一个元素（即只含 $0$），$R=\{0\}$ 的环。
    \item \textbf{域 (Field)}: 设 $R$ 为具有至少两个元的外交换环。集合 $R^* = R \setminus \{0\}$，如果满足 $(R^*, \cdot)$ 构成一个阿贝尔群，则 $R$ 称为一个域，通常记为 $F$。
\end{enumerate}
\end{cndefinition}
直观理解：在域中，除了零以外的所有元素都可以做除法。域是代数中最完美的结构之一。有理数 $\mathbb{Q}$ 是域，整数 $\mathbb{Z}$ 是整环但不是域。

\begin{cnexample}
$\mathbb{Z}_p$（$p$ 为素数）关于加法乘法都构成阿贝尔群，为域。$\mathbb{Q}$ 和 $\mathbb{Q}[i]$ 在普通的加法和乘法下均为域。
\end{cnexample}

\section{环同态 (Ring Homomorphism)}

\begin{cndefinition}[环同态]
设 $R$ 和 $S$ 是两个环。映射 $f: R \to S$ 称为一个\textbf{环同态}，如果对于任意 $\alpha, \beta \in R$，都满足：
\begin{enumerate}
    \item $f(\alpha + \beta) = f(\alpha) + f(\beta)$ （保持加法运算）
    \item $f(\alpha\beta) = f(\alpha)f(\beta)$ （保持乘法运算）
\end{enumerate}
\end{cndefinition}

\begin{cnproperty}[环同态的基本性质]
设 $f: R \to S$ 是环同态，则：
\begin{enumerate}
    \item $f(0_R) = 0_S$，且 $f(-\alpha) = -f(\alpha)$。
    \item 如果 $f$ 分别是单射、满射或双射，则 $f$ 分别称为单同态、满同态或同构。如果是自映射 $f:R \to R$ 则称\textbf{自同态}，双射自映射称\textbf{自同构}。
    \item 单同态 $f$ 有时被称为由 $R$ 到 $S$ 的\textbf{嵌入} (embedding)。
    \item 像集 $f(R)$ 是 $S$ 的子环。
    \item 如果 $R$ 和 $S$ 都是含幺环，且 $f$ 是\textbf{满同态}，则必有 $f(1_R) = 1_S$。进一步地，若 $u \in U(R)$，有 $f(u^{-1}) = f(u)^{-1}$。
\end{enumerate}
\end{cnproperty}
\begin{proof}
\begin{enumerate}
    \item 结合群同态性质，$f(0_R) = f(0_R + 0_R) = f(0_R) + f(0_R)$。在 $S$ 的加法群中抵消 $f(0_R)$，即得 $f(0_R) = 0_S$。又 $0_S = f(0_R) = f(\alpha + (-\alpha)) = f(\alpha) + f(-\alpha)$，即 $f(-\alpha)$ 是 $f(\alpha)$ 的加法逆元。
    \item 满足环 $R$ 含幺且 $f$ 为满射。对任意 $s \in S$，存在 $r \in R$ 使 $f(r) = s$。
    $$s = f(r) = f(r \cdot 1_R) = f(r)f(1_R) = s f(1_R)$$
    同时 $s = f(1_R \cdot r) = f(1_R)f(r) = f(1_R) s$。可见 $f(1_R)$ 就是 $S$ 的幺元 $1_S$。对于单位元 $u$，有 $1_S = f(1_R) = f(uu^{-1}) = f(u)f(u^{-1})$，因此 $f(u^{-1}) = f(u)^{-1}$。
\end{enumerate}
\end{proof}

\section{子环与特征 (Characteristic)}

\begin{cndefinition}[子环与子域]
设 $S$ 是环 $R$ 的子集。如果在 $R$ 的加法和乘法下，$(S, +, \cdot)$ 本身也构成一个环，则称 $S$ 是 $R$ 的\textbf{子环}。相似地，如果域 $F$ 的子集 $S$ 自己本身按 $F$ 上的运算构成一个域，则称 $S$ 为 $F$ 的\textbf{子域}。此时称 $F$ 为 $S$ 的\textbf{扩域} (Extension field)。
\end{cndefinition}
\begin{cnexample}
$\mathbb{Z}$ 是 $\mathbb{Q}$ 的子环；$\mathbb{Z}$ 也是高斯整数环 $\mathbb{Z}[i]$ 的子环。对于域，$\mathbb{Q} \subset \mathbb{Q}[i] \subset \mathbb{C}$ 以及 $\mathbb{Q} \subset \mathbb{R} \subset \mathbb{C}$ 构成子域与扩域链。
\end{cnexample}

\begin{cndefinition}[环的特征]
对于环 $R$，如果存在一个最小的正整数 $m > 0$，使得对任意 $r \in R$，都有 $mr = 0$（即 $r$ 自加 $m$ 次等于零元），则称 $m$ 为环 $R$ 的\textbf{特征 (Characteristic)}，记作 $\text{char}(R) = m$。
换言之，$m$ 是所有元素加法阶的最小公倍数。显然若 $R$ 为含幺环，则等价于存在最小的 $m>0$ 使得 $m \cdot 1 = 0$。如果这样的 $m$ 不存在，则规定 $\text{char}(R) = 0$。
\end{cndefinition}
\begin{cnexample}
$\text{char}(\mathbb{Z}) = 0$，$\text{char}(\mathbb{Z}_n) = n$。
\end{cnexample}

\begin{cntheorem}
如果 $R$ 是一个整环，且 $\text{char}(R) \neq 0$，则 $\text{char}(R)$ 必须是一个素数。
\end{cntheorem}
\begin{proof}
设 $\text{char}(R) = m > 0$。如果 $m$ 不是素数，则可以分解为 $m = a \cdot b$，其中 $1 < a < m$ 且 $1 < b < m$。
由于特征的定义 $m \cdot 1 = 0$，我们可以写：
$$(a \cdot 1)(b \cdot 1) = (ab) \cdot 1 = m \cdot 1 = 0$$
但 $R$ 是整环，没有零因子，所以必须有 $(a \cdot 1) = 0$ 或 $(b \cdot 1) = 0$。如果 $a \cdot 1 = 0$，显然对任何 $r \in R$ 有 $ar = a \cdot (1 \cdot r) = (a \cdot 1)r = 0 \cdot r = 0$。这与 $m$ 是满足该性质的“最小”正整数相矛盾。因此 $m$ 不能是合数，$m$ 必定是素数。
\end{proof}

\begin{cncorollary}
因为域必定是整环，所以任何域的特征要么是 $0$，要么是某个素数 $p$。
\end{cncorollary}

\begin{cnproperty}[Freshman's Dream / 新生之梦]
如果 $R$ 是一个交换环，且 $\text{char}(R) = p$ （$p$ 为素数）。那么对于任意 $\alpha, \beta \in R$，恒有公式
$$(\alpha + \beta)^p = \alpha^p + \beta^p$$
\end{cnproperty}
\begin{proof}
由二项式定理（在交换环中适用），我们有：
$$(\alpha + \beta)^p = \sum_{k=0}^{p} \binom{p}{k} \alpha^{p-k} \beta^k$$
对于 $0 < k < p$，组合数 $\binom{p}{k} = \frac{p!}{k!(p-k)!}$ 一定是 $p$ 的倍数，因为分子包含素数因子 $p$，而分母的任何项都消不掉 $p$。由于环 $R$ 的特征是 $p$，即意味着对于任意 $r \in R$，有 $pr = 0$。因此中间项所有系数 $\binom{p}{k}$ 乘以任意元素均等于 $0$。化简后只存留第一项和最后一项，即 $\alpha^p + \beta^p$。
\end{proof}

\end{document}
